% ---------------------------------------------------------------------
\chapter{Discusiones y conclusiones}\label{cap:conclusiones}

\section{Conclusiones}

En este tipo de problema, la preparación de los datos es la que decide el resultado: vemos cómo unas pocas variables clínicas que estén bien construidas, usando un modelo lineal, son capaces de alcanzar un AUROC de entre \num{0.91} y \num{0.93}.

Por objetivos, hemos concluido que:

\begin{enumerate}[label=\textbf{O\arabic*}, leftmargin=2.4em]
  \item \textbf{ETL y calidad.} El conjunto de datos es de buena calidad: no contiene ningún valor ausente, tanto el ruido como la deriva son pequeños y todos los registros nos permiten construir un latido medio fiable. Durante el control de calidad fuimos capaces de detectar un registro duplicado, que podría haber causado una fuga entre el entrenamiento y el test.
  \item \textbf{Preprocesado.} El filtro elegido apenas influye, mientras que el detector de latidos sí que lo hace. En una señal muestreada a \SI{100}{\hertz}, la precisión temporal de la segmentación importa más que cómo de limpia esté la señal.
  \item \textbf{Selección de instancias.} Solo se descarta un \SI{0.9}{\percent} de los latidos, sin que exista un sesgo entre clases, y ninguno de los ocho criterios evaluados cambia el resultado. La mediana ya actúa como un estimador robusto, y la selección se mantiene como salvaguarda y por trazabilidad.
  \item \textbf{Características.} Sacamos en claro que la pendiente del segmento ST y la duración del QRS son capaces de discriminar mucho más que la elevación absoluta del ST, lo que coincide con los criterios de forma de la literatura clínica, como el índice de Corrado \cite{zorzi2015}. Los cuatro métodos de selección comparados rinden prácticamente igual, y coinciden en que con apenas 5 variables se conserva casi el \SI{99}{\percent} del rendimiento, así que bastaría con usar uno sencillo y barato, como puede ser el ANOVA.
  \item \textbf{Visualización.} Al final, ninguna de las proyecciones logró separar las dos clases. La PCA muestra que las direcciones de máxima varianza no son las discriminantes, y t-SNE y UMAP solo nos dejan ver agrupaciones locales parciales.
  \item \textbf{Modelo de referencia.} El LDA con contracción alcanza un AUROC de \num{0.928} (IC 95\,\%: \numrange{0.88}{0.96}), un valor muy superior al obtenido con la puntuación basada en los criterios clínicos (\num{0.75}).
\end{enumerate}

\section{Limitaciones}

\begin{itemize}
  \item Al trabajar a \SI{100}{\hertz}, cada muestra equivale a \SI{10}{\milli\second}. Por eso la duración del QRS solo puede tomar valores escalonados, y el punto J se localiza con una incertidumbre de una muestra.
  \item Los datos que tenemos son solo de un centro y de un ECG por paciente, y no tenemos ningún tipo de validación externa de estos registros.
  \item El tamaño muestral es muy bajo también: contando solo con unos 15 positivos por \eng{fold} de test, las diferencias de AUROC que sean menores de \numrange{0.02}{0.03} no llegan a ser concluyentes.
\end{itemize}

\section{Continuación con el Trabajo de Fin de Grado}

El proceso y el conjunto de datos construidos aquí representan la base de la parte comparativa de mi TFG. Con el mismo protocolo de evaluación, en el TFG compararé siete familias de métodos: reglas clínicas, modelos clásicos, autoencoders, redes convolucionales y de atención sobre la señal, el modelo preentrenado ECG-FM \cite{mckeen2025}, redes sobre imágenes del latido y modelos multimodales de visión-lenguaje \cite{costa2026vlm}. También se analizarán la interpretabilidad de los modelos, la calibración de sus probabilidades y los pacientes que todos los métodos clasifican mal. El valor de referencia que fijamos en el \cref{cap:modelo} es la cifra que tratarán de superar todos estos nuevos y más complejos modelos.
